\section*{Schedules}
\label{sec:schedules}

\scheduleSection{Limits and convergence}{6}{Sequences and series in \(\RR\) and \(\CC\). Sums, products and quotients. Absolute convergence; absolute convergence implies convergence. The Bolzano--Weierstrass theorem and applications (the General Principle of Convergence). Comparison and ratio tests, alternating series test.}

\scheduleSection{Continuity}{3}{Continuity of real- and complex-valued functions defined on subsets of \(\RR\) and \(\CC\). The intermediate value theorem. A continuous function on a closed bounded interval is bounded and attains its bounds.}

\scheduleSection{Differentiability}{5}{Differentiability of functions from \(\RR\) to \(\RR\). Derivative of sums and products. The chain rule. Derivative of the inverse function. Rolle's theorem; the mean value theorem. One-dimensional version of the inverse function theorem. Taylor's theorem from \(\RR\) to \(\RR\); Lagrange's form of the remainder. Complex differentiation.}

\scheduleSection{Power series}{4}{Complex power series and radius of convergence. Exponential, trigonometric and hyperbolic functions, and relations between them. *Direct proof of the differentiability of a power series within its circle of convergence*.}

\scheduleSection{Integration}{6}{Definition and basic properties of the Riemann integral. A non-integrable function. Integrability of monotonic functions. Integrability of piecewise-continuous functions. The fundamental theorem of calculus. Differentiation of indefinite integrals. Integration by parts. The integral form of the remainder in Taylor's theorem. Improper integrals.}